% GENERATED by scripts/build-lecture-slides.py; do not edit. Portable profile: XeTeX, theme and fonts from the core.
\makeatletter\IfFileExists{core/beamer/beamerthememx.sty}{\def\mxcorepath{core/}\def\input@path{{core/beamer/}{core/}}}{\def\mxcorepath{../core/}\def\input@path{{../core/beamer/}{../core/}}}\makeatother
\documentclass[aspectratio=169,10pt,t]{beamer}
\usetheme{mx}
\usepackage{graphicx,booktabs}
\usepackage[french]{babel}
\setlength{\parskip}{.45em}
\hypersetup{pdfauthor={},pdfsubject={Lecture slides},pdfkeywords={},pdftitle={Calcul vectoriel et produit scalaire}}
\providecommand{\N}{\mathbb{N}}
\providecommand{\Z}{\mathbb{Z}}
\providecommand{\D}{\mathbb{D}}
\providecommand{\Q}{\mathbb{Q}}
\providecommand{\R}{\mathbb{R}}
\providecommand{\C}{\mathbb{C}}
\providecommand{\vect}[1]{\overrightarrow{#1}}
\providecommand{\vi}{\vec{\imath}}
\providecommand{\vj}{\vec{\jmath}}
\providecommand{\vk}{\vec{k}}
\providecommand{\norm}[1]{\left\lVert #1\right\rVert}
\providecommand{\vc}[2]{\begin{pmatrix}#1\\#2\end{pmatrix}}
\providecommand{\deter}[2]{\det\left(#1,#2\right)}
\providecommand{\abs}[1]{\left\lvert #1\right\rvert}
\providecommand{\dd}{\,\mathrm{d}}
\providecommand{\e}{\mathrm{e}}
\providecommand{\Card}{\operatorname{Card}}
\providecommand{\E}{\operatorname{E}}
\providecommand{\V}{\operatorname{V}}
\providecommand{\degre}{^{\circ}}
\providecommand{\un}[1]{\,\mathrm{#1}}
\providecommand{\sci}[2]{#1\times10^{#2}}
\providecommand{\pgcd}{\operatorname{PGCD}}
\providecommand{\ppcm}{\operatorname{PPCM}}
\title{Calcul vectoriel et produit scalaire}
\subtitle{Cours et exemples}
\date{}
\author{}
\begin{document}
\mxtitleframe
\begin{frame}{Objectifs}
Décomposer un vecteur dans une base et utiliser les expressions du produit scalaire pour démontrer des relations métriques.
\end{frame}
\begin{frame}{Combinaisons linéaires et bases}
Une combinaison linéaire de \(\overrightarrow{u},\overrightarrow{v}\) est \(a\overrightarrow{u}+b\overrightarrow{v}\) avec \(a,b\in\mathbb{R}\). Deux vecteurs non colinéaires forment une base du plan : chaque vecteur s’écrit d’une manière unique \(a\overrightarrow{u}+b\overrightarrow{v}\).
\end{frame}
\begin{frame}{Combinaisons linéaires et bases}
Dans un repère, \(\overrightarrow{u}=(x;y)\) et \(\overrightarrow{v}=(x';y')\) sont non colinéaires si et seulement si \(xy'-yx'\ne0\).
\end{frame}
\begin{frame}{Combinaisons linéaires et bases}
Le système \(a\overrightarrow{u}+b\overrightarrow{v}=\overrightarrow{w}\) a deux équations. L’élimination fournit une solution unique si le déterminant \(xy'-yx'\) est non nul. S’il est nul, les colonnes sont proportionnelles (ou l’un des vecteurs est nul), donc colinéaires.
\end{frame}
\begin{frame}{Combinaisons linéaires et bases}
\(\overrightarrow{u}=(1;1)\), \(\overrightarrow{v}=(1;-1)\) forment une base. Pour \(\overrightarrow{w}=(3;1)\), \(a+b=3\) et \(a-b=1\), donc \(\overrightarrow{w}=2\overrightarrow{u}+\overrightarrow{v}\). Cette base n’est pas un repère orthonormé : ses deux vecteurs ont norme \(\sqrt2\).
\end{frame}
\begin{frame}{Quatre expressions d’un même produit}
Pour deux vecteurs non nuls, \(\overrightarrow{u}\cdot\overrightarrow{v}=\left\| \overrightarrow{u} \right\|\left\| \overrightarrow{v} \right\|\cos\theta\), où \(\theta\) est l’angle entre eux. Si l’un est nul, le produit scalaire est nul.
\end{frame}
\begin{frame}{Quatre expressions d’un même produit}
Dans un repère orthonormé, \(\overrightarrow{u}\cdot\overrightarrow{v}=xx'+yy'\). Pour \(O\ne A\), avec \(H\) projeté orthogonal de \(B\) sur la droite \((OA)\) et \(\overrightarrow{u}=\overrightarrow{OA},\overrightarrow{v}=\overrightarrow{OB}\), le produit vaut \(OA\) multiplié par la longueur algébrique \(OH\) orientée vers \(A\). Enfin,
\[
\overrightarrow{u}\cdot\overrightarrow{v}=\dfrac12\big(\left\| \overrightarrow{u}+\overrightarrow{v} \right\|^2-\left\| \overrightarrow{u} \right\|^2-\left\| \overrightarrow{v} \right\|^2\big).
\]
Le produit est symétrique et bilinéaire; \(\overrightarrow{u}\cdot\overrightarrow{u}=\left\| \overrightarrow{u} \right\|^2\); deux vecteurs sont orthogonaux si et seulement si leur produit est nul (y compris le vecteur nul).
\end{frame}
\begin{frame}{Quatre expressions d’un même produit}
Décomposer dans deux directions orthonormées donne \(xx'+yy'\). Développer \((x+x')^2+(y+y')^2\) donne l’identité par les normes. La projection donne \(OH=OB\cos\theta\) avec son signe, donc la même expression géométrique. La bilinéarité résulte de la formule des coordonnées.
\end{frame}
\begin{frame}{Relations métriques dans un triangle}
Dans un triangle \(ABC\),
\[
BC^2=AB^2+AC^2-2AB\,AC\cos\widehat{BAC}.
\]
Si \(I\) est milieu de \([BC]\), alors \(AB^2+AC^2=2AI^2+BC^2/2\).
\end{frame}
\begin{frame}{Relations métriques dans un triangle}
Al-Kashi résulte du carré de \(\overrightarrow{BC}=\overrightarrow{AC}-\overrightarrow{AB}\). Pour la médiane, \(\overrightarrow{AB}=\overrightarrow{AI}+\overrightarrow{IB}\) et \(\overrightarrow{AC}=\overrightarrow{AI}-\overrightarrow{IB}\); en additionnant les carrés, les produits croisés s’annulent.
\end{frame}
\begin{frame}{Relations métriques dans un triangle}
Pour un triangle non aplati, \(\mathcal A=\dfrac12 AB\,AC\sin\widehat{BAC}\) et
\[
\dfrac{BC}{\sin\widehat{BAC}}=
\dfrac{CA}{\sin\widehat{ABC}}=
\dfrac{AB}{\sin\widehat{ACB}}.
\]
\end{frame}
\begin{frame}{Relations métriques dans un triangle}
La hauteur relative à \(AB\) vaut \(AC\sin\widehat{BAC}\). Écrire l’aire avec chacun des trois angles puis diviser par les produits des côtés (tous non nuls) donne la relation des sinus.
\end{frame}
\begin{frame}{Relations métriques dans un triangle}
Avec des coordonnées dans un repère orthonormé, calculer d’abord les vecteurs puis leur produit. Pour un angle, diviser par les deux normes non nulles. Pour projeter \(C\) sur la droite \((AB)\), \(A\ne B\), poser \(H=A+t\overrightarrow{AB}\); la condition \(\overrightarrow{HC}\cdot\overrightarrow{AB}=0\) impose
\[
t=\dfrac{\overrightarrow{AC}\cdot\overrightarrow{AB}}{AB^2}.
\]
Si \(t\notin[0;1]\), le pied appartient à la droite mais pas au segment; la hauteur d’un triangle peut être extérieure.
\end{frame}
\begin{frame}{Relations métriques dans un triangle}
Pour \(A(1;1)\), \(B(5;3)\), \(C(2;4)\), \(\overrightarrow{AB}=(4;2)\), \(\overrightarrow{AC}=(1;3)\). Le produit vaut \(10\) et \(AB^2=20\), donc \(t=1/2\), \(H(3;2)\). On vérifie \(\overrightarrow{HC}=(-1;2)\) et son produit avec \((4;2)\) vaut \(0\). La distance \(CH=\sqrt5\) est la hauteur relative à \(AB=2\sqrt5\), d’où l’aire \(5\) unités d’aire. Le même produit donne \(\cos\widehat{BAC}=10/(\sqrt{20}\sqrt{10})=1/\sqrt2\) : l’angle vaut \(\pi/4\).
\end{frame}
\begin{frame}{Relations métriques dans un triangle}
Les coordonnées d’un vecteur sont les coefficients de sa décomposition dans la base choisie. Si la base n’est pas orthonormée, les produits des vecteurs de base interviennent. On peut revenir aux coordonnées dans un repère orthonormé pour éviter une formule mal appliquée. Une base sert à représenter les vecteurs; elle ne garantit pas des unités de longueur égales.
\end{frame}
\end{document}
